\chapter{Self-shrinkers in higher codimension}\label{sec:examples-higher-codimension-shrinkers}
These examples use the vector convention $\vec H=\operatorname{tr}B=\Delta X=-X^\perp/2$ for a smooth Euclidean immersion $X:M^n\to\mathbb R^N$, where $n\geq1$. Thus the squared curvature is $|B|^2=\sum_{i,j}|B(e_i,e_j)|^2$. The primary models are \cite{CL12}, Section 5, pp.~8--9, rescaled from that paper's convention $\vec H=-X^\perp$. The scalar hypersurface formulas of Section \ref{sec:shrinker-curvature-gap} cannot simply be used for a higher-rank normal bundle.

\begin{proposition}[Minimal submanifolds of a sphere give shrinkers]\label{prop:spherical-minimal-shrinker-equivalence}
Let $X:M^n\to S^{N-1}(R)\subset\mathbb R^N$, $n\geq1$ and $R>0$, be a smooth immersion. It is a Euclidean self-shrinker if and only if $R=\sqrt{2n}$ and it is minimal in that sphere. For this radius,
\[
 |B^{\mathbb R^N}|^2=|B^{S^{N-1}(R)}|^2+\frac12.
\]
\end{proposition}
\begin{proof}
The sphere's vector second fundamental form in Euclidean space is $-(1/R^2)gX$. Decomposing the ambient derivative first relative to $M$ in the sphere and then relative to the sphere in Euclidean space gives
\[
 B^{\mathbb R^N}_{ij}=B^{S^{N-1}(R)}_{ij}-\frac{\delta_{ij}}{R^2}X,
 \qquad \vec H^{\mathbb R^N}=\vec H^{S^{N-1}(R)}-\frac n{R^2}X.
\]
The two normal summands are perpendicular. Also $X^\perp=X$, since $|X|$ is constant. The equation $\vec H^{\mathbb R^N}=-X/2$ therefore says separately that the spherical mean curvature vanishes and $n/R^2=1/2$. Squaring the orthogonal decomposition of $B$ gives the norm identity, since the radial summand has squared norm $n/R^2=1/2$.
\end{proof}

\begin{example}[Products of round factors and a Euclidean factor]\label{ex:product-sphere-shrinkers}
Let $p\geq1$, $d_1,\ldots,d_p\geq1$ and $q\geq0$ be integers, and put $m=\sum_i d_i$, $n=m+q$. The embedded product
\[
 M=\prod_{i=1}^p S^{d_i}(\sqrt{2d_i})\times\mathbb R^q\subset\mathbb R^{n+p}
\]
is a complete proper self-shrinker. Its curvature satisfies $|B|^2=p/2$ and $|\vec H|^2=m/2$. It has polynomial area growth, is compact exactly when $q=0$, and in that case is minimal in $S^{n+p-1}(\sqrt{2n})$.
\end{example}
\begin{proof}
For arbitrary radii $r_i$, the mean-curvature vector has $i$th spherical block $-d_iX_i/r_i^2$ and zero Euclidean block. The normal projection of $X=(X_1,\ldots,X_p,y)$ is $(X_1,\ldots,X_p,0)$. Thus the shrinker equation holds exactly when $r_i^2=2d_i$ for every $i$. The different radial normals are orthogonal, so
\[
 |B|^2=\sum_i\frac{d_i}{r_i^2}=p/2,
 \qquad |\vec H|^2=\sum_i\frac{d_i^2}{r_i^2}=m/2.
\]
A product of complete round metrics and a Euclidean metric is complete. Compactness of the spherical factors and boundedness of the Euclidean coordinate on ambient compact sets prove properness. The portion in an ambient ball is bounded by the product of the total spherical volumes and a radius-$R$ Euclidean ball volume, hence has growth at most $C(1+R)^q$. When $q=0$, the squared ambient radius is $\sum_i2d_i=2n$ and Proposition \ref{prop:spherical-minimal-shrinker-equivalence} proves spherical minimality. Compare \cite{CL12}, Examples 5.1--5.2.
\end{proof}

\begin{proposition}[Gaussian area of the product models]\label{prop:product-shrinker-gaussian-area}
Write $\sigma_d=\mathcal H^d(S^d(1))$. The normalized Gaussian area of Example \ref{ex:product-sphere-shrinkers} is
\[
 F_{0,1}(M)=\prod_{i=1}^p\left[(4\pi)^{-d_i/2}e^{-d_i/2}\sigma_{d_i}(2d_i)^{d_i/2}\right].
\]
It is independent of the dimension $q$ of the Euclidean factor.
\end{proposition}
\begin{proof}
The product area element factors, and $|X|^2=2\sum_i d_i+|y|^2$. Tonelli's theorem separates the positive Gaussian integral into the spherical factors and
\[
 (4\pi)^{-q/2}\int_{\mathbb R^q}e^{-|y|^2/4}\,dy=1.
\]
The radius-$\sqrt{2d_i}$ sphere has area $\sigma_{d_i}(2d_i)^{d_i/2}$, on which the Gaussian exponential is the constant $e^{-d_i/2}$. Multiplying gives the formula. The statement computes Gaussian area at center zero and scale one; it does not invoke a separate theorem identifying an entropy supremum over all centers and scales.
\end{proof}

\begin{example}[The Clifford torus shrinker in four dimensions]\label{ex:clifford-torus-shrinker}
The map
\[
 X(\theta,\varphi)=(\sqrt2\cos\theta,\sqrt2\sin\theta,
 \sqrt2\cos\varphi,\sqrt2\sin\varphi)
\]
on $(\mathbb R/2\pi\mathbb Z)^2$ is a flat embedded torus in $\mathbb R^4$ with $|B|^2=1$, $|\vec H|^2=1$, area $8\pi^2$ and $F_{0,1}=2\pi/e$. It is minimal in $S^3(2)$, where its squared second fundamental form is $1/2$.
\end{example}
\begin{proof}
The metric is $2(d\theta^2+d\varphi^2)$, so it is flat with area $8\pi^2$. Apply Example \ref{ex:product-sphere-shrinkers} with two circle factors. Since $|X|^2=4$, direct substitution gives $F_{0,1}=(4\pi)^{-1}e^{-1}8\pi^2=2\pi/e$. Proposition \ref{prop:spherical-minimal-shrinker-equivalence} subtracts $1/2$ from the Euclidean curvature norm to give the spherical value. This is a codimension-two shrinker and its Euclidean curvature exceeds the hypersurface gap threshold.
\end{proof}

\begin{proposition}[The compact mean-curvature threshold]\label{prop:compact-shrinker-mean-curvature-threshold}
On every connected closed Euclidean self-shrinker $M^n$, in any codimension,
\[
 \frac1{\operatorname{Vol}(M)}\int_M|\vec H|^2\,d\mu=\frac n2.
\]
If $|\vec H|^2$ is everywhere at most $n/2$, everywhere at least $n/2$, or constant, then the image lies in $S^{N-1}(\sqrt{2n})$ and is minimal there. See \cite{CL12}, Proposition 5.2, for the source normalization.
\end{proposition}
\begin{proof}
The immersion identities give
\[
 \Delta|X|^2=2n+2\langle X,\vec H\rangle=2n-4|\vec H|^2.
\]
Integrating on a closed manifold gives the average identity. Each stated pointwise condition forces $|\vec H|^2=n/2$ everywhere, so $|X|^2$ is harmonic and hence constant on the connected compact manifold. It follows that $X^T=\nabla|X|^2/2=0$. The shrinker equation now gives $|X|^2=4|\vec H|^2=2n$, and Proposition \ref{prop:spherical-minimal-shrinker-equivalence} gives spherical minimality. This conclusion does not say that an arbitrary-codimension example is a round hypersphere.
\end{proof}

\begin{proposition}[There is no closed Euclidean self-expander]\label{prop:no-closed-higher-codimension-expander}
For $n\geq1$, no nonempty closed immersed $n$-manifold in Euclidean space satisfies $\vec H=X^\perp/2$, in any codimension.
\end{proposition}
\begin{proof}
For an expander the same immersion calculation gives $\Delta|X|^2=2n+|X^\perp|^2>0$. Its integral on a closed manifold would be positive, whereas the divergence theorem makes that integral zero. This proves the contradiction and the observation in \cite{CL12}, Proposition 5.3, directly with the present normalization.
\end{proof}

\begin{lemma}[The Gaussian radial identity in any codimension]\label{lem:higher-codimension-shrinker-radial-moment}
Let $X:M^n\to\mathbb R^N$ be a connected complete smooth boundaryless self-shrinker with polynomial area growth, counting area on the parameter manifold. With $\rho=(4\pi)^{-n/2}e^{-|X|^2/4}$,
\[
 \int_M |X|^2\rho\,d\mu=2n\int_M\rho\,d\mu.
\]
Neither a global unit normal nor hypersurface codimension is required.
\end{lemma}
\begin{proof}
The identities $\Delta X=\vec H=-X^\perp/2$ and $|DX|^2=n$ give $\Delta|X|^2=2n-|X^\perp|^2$. Since $\nabla|X|^2=2X^T$, subtraction of the Gaussian drift gives
\[
 \mathcal D|X|^2=2n-|X|^2,\qquad
 \mathcal D=\Delta-\tfrac12\nabla_{X^T}
 =\rho^{-1}\operatorname{div}(\rho\nabla\,\cdot\,).
\]
Polynomial area growth implies integrability of $(1+|X|^2)\rho$ by exactly the annular estimate in Lemma \ref{lem:gaussian-polynomial-moments}; that estimate uses only area, not codimension. Take intrinsic compact cutoffs $\eta_R$ equal to one on a radius-$R$ ball, zero outside radius $2R$, with $|\nabla\eta_R|\leq2/R$. Completeness supplies compact support. Weighted integration by parts gives
\[
 \int\eta_R(2n-|X|^2)\rho
 =-\int\langle\nabla\eta_R,2X^T\rangle\rho.
\]
The right side has absolute value at most $(4/R)\int|X|\rho$, tending to zero. Dominated convergence gives the conclusion on the left. This explicitly extends the hypersurface computation instead of invoking a hypersurface statement for a higher-codimension immersion.
\end{proof}

\begin{proposition}[A complete mean-curvature bound forces spherical containment]\label{prop:complete-shrinker-mean-curvature-threshold}
Under the hypotheses of Lemma \ref{lem:higher-codimension-shrinker-radial-moment}, assume $|\vec H|^2\geq n/2$ everywhere. Then $|\vec H|^2=n/2$, $X^T=0$, and the immersed image is a minimal submanifold of $S^{N-1}(\sqrt{2n})$. This is the rescaled assertion of \cite{CL12}, Proposition 5.1. It does not classify all minimal submanifolds of that sphere.
\end{proposition}
\begin{proof}
The shrinker equation gives the pointwise identity
\[
 |X|^2=|X^T|^2+4|\vec H|^2\geq2n.
\]
The preceding lemma says that its Gaussian average is exactly $2n$. The nonnegative continuous function $|X|^2-2n$ therefore has zero weighted integral and vanishes everywhere. Both nonnegative terms in $|X^T|^2+4(|\vec H|^2-n/2)=0$ vanish. The fixed-radius criterion in Proposition \ref{prop:spherical-minimal-shrinker-equivalence} now gives spherical minimality. The Euclidean-factor models of Example \ref{ex:product-sphere-shrinkers}, with $q>0$, have $|\vec H|^2=(n-q)/2<n/2$ and unbounded image, showing why that lower threshold cannot simply be omitted.
\end{proof}
